\documentclass[12pt, a4paper]{article}

%% ─── Packages ────────────────────────────────────────────────────────────────
\usepackage[T1]{fontenc}
\usepackage[utf8]{inputenc}
\usepackage{lmodern}
\usepackage{microtype}
\usepackage{amsmath, amssymb, amsthm}
\usepackage{mathtools}
\usepackage{geometry}
\usepackage{hyperref}
\usepackage{xcolor}
\usepackage{listings}
\usepackage{booktabs}
\usepackage{array}
\usepackage{enumitem}
\usepackage{cite}
\usepackage{graphicx}
\usepackage{tikz}
\usepackage{setspace}

\geometry{
  a4paper,
  top=2.5cm, bottom=2.5cm,
  left=2.8cm, right=2.8cm
}

\hypersetup{
  colorlinks = true,
  linkcolor  = blue!70!black,
  citecolor  = blue!70!black,
  urlcolor   = blue!70!black
}

%% ─── Lean listing style ──────────────────────────────────────────────────────
\definecolor{leanblue}{RGB}{0, 0, 180}
\definecolor{leanorange}{RGB}{200, 100, 0}
\definecolor{leangreen}{RGB}{0, 130, 0}
\definecolor{leangray}{RGB}{120, 120, 120}
\definecolor{codebg}{RGB}{248, 248, 252}

\lstdefinelanguage{Lean4}{
  keywords={def, theorem, lemma, axiom, noncomputable, import, open,
            section, namespace, end, where, by, exact, intro, apply,
            rfl, simp, linarith, ring, constructor, rcases, obtain,
            have, show, let, fun, if, then, else, structure, inductive,
            instance, class, universe, variable, set_option, attribute,
            scoped, notation, abbrev, example, calc, match, with,
            return, do, for, unless, at, from, in, forall, exists,
            True, False, Prop, Type, Sort},
  keywordstyle=\color{leanblue}\bfseries,
  comment=[l]{--},
  commentstyle=\color{leangray}\itshape,
  stringstyle=\color{leangreen},
  sensitive=true,
  morecomment=[s]{/-}{-/},
  basicstyle=\ttfamily\small,
  backgroundcolor=\color{codebg},
  frame=single,
  framerule=0.4pt,
  rulecolor=\color{gray!40},
  xleftmargin=8pt,
  xrightmargin=4pt,
  framexleftmargin=6pt,
  numbers=none,
  keepspaces=true,
  breaklines=true,
  breakatwhitespace=true,
  tabsize=2,
  showstringspaces=false,
  inputencoding=utf8,
  extendedchars=true,
  literate=
    {∀}{{$\forall$}}1
    {∃}{{$\exists$}}1
    {∈}{{$\in$}}1
    {∉}{{$\notin$}}1
    {⊆}{{$\subseteq$}}1
    {∧}{{$\wedge$}}1
    {∨}{{$\vee$}}1
    {¬}{{$\neg$}}1
    {≠}{{$\neq$}}1
    {≤}{{$\leq$}}1
    {≥}{{$\geq$}}1
    {≡}{{$\equiv$}}1
    {≈}{{$\approx$}}1
    {→}{{$\to$}}1
    {←}{{$\leftarrow$}}1
    {↦}{{$\mapsto$}}1
    {⟨}{{$\langle$}}1
    {⟩}{{$\rangle$}}1
    {⊤}{{$\top$}}1
    {⊥}{{$\bot$}}1
    {ℝ}{{$\mathbb{R}$}}1
    {ℂ}{{$\mathbb{C}$}}1
    {ℕ}{{$\mathbb{N}$}}1
    {ℤ}{{$\mathbb{Z}$}}1
    {ℚ}{{$\mathbb{Q}$}}1
    {Δ}{{$\Delta$}}1
    {δ}{{$\delta$}}1
    {Λ}{{$\Lambda$}}1
    {λ}{{$\lambda$}}1
    {Ξ}{{$\Xi$}}1
    {ξ}{{$\xi$}}1
    {Ω}{{$\Omega$}}1
    {ω}{{$\omega$}}1
    {Σ}{{$\Sigma$}}1
    {σ}{{$\sigma$}}1
    {Π}{{$\Pi$}}1
    {π}{{$\pi$}}1
    {ε}{{$\varepsilon$}}1
    {α}{{$\alpha$}}1
    {β}{{$\beta$}}1
    {γ}{{$\gamma$}}1
    {ρ}{{$\rho$}}1
    {τ}{{$\tau$}}1
    {ψ}{{$\psi$}}1
    {φ}{{$\varphi$}}1
    {μ}{{$\mu$}}1
    {ν}{{$\nu$}}1
    {ζ}{{$\zeta$}}1
    {η}{{$\eta$}}1
    {θ}{{$\theta$}}1
    {ℒ}{{$\mathcal{L}$}}1
    {²}{{$^2$}}1
    {³}{{$^3$}}1
    {∞}{{$\infty$}}1
    {§}{{\S}}1
    {·}{{$\cdot$}}1
    {¶}{{\P}}1
    {★}{{$\star$}}1
    {ᶜ}{{$^{\mathsf{c}}$}}1
    {ᵇ}{{$_b$}}1
    {á}{{\'a}}1
    {é}{{\'e}}1
    {í}{{\'i}}1
    {ó}{{\'o}}1
    {ú}{{\'u}}1
}

%% ─── Theorem environments ────────────────────────────────────────────────────
\theoremstyle{plain}
\newtheorem{theorem}{Theorem}[section]
\newtheorem{lemma}[theorem]{Lemma}
\newtheorem{corollary}[theorem]{Corollary}
\newtheorem{proposition}[theorem]{Proposition}

\theoremstyle{definition}
\newtheorem{definition}[theorem]{Definition}
\newtheorem{example}[theorem]{Example}
\newtheorem{remark}[theorem]{Remark}

%% ─── Macros ──────────────────────────────────────────────────────────────────
\newcommand{\N}{\mathcal{N}}
\newcommand{\R}{\mathbb{R}}
\newcommand{\C}{\mathbb{C}}
\newcommand{\Z}{\mathbb{Z}}
\newcommand{\re}{\operatorname{Re}}
\newcommand{\zetaord}{\operatorname{ord}}
\newcommand{\lean}[1]{\texttt{#1}}
\newcommand{\proved}{\textbf{[Proved]}}
\newcommand{\axiom}{\textbf{[Axiom]}}
\newcommand{\Mathlib}{\textsf{Mathlib}}
\newcommand{\Lean}{\textsf{Lean~4}}
\newcommand{\Aristotle}{\textsf{Aristotle}}
\newcommand{\ShiftedSpec}{\mathcal{S}_{1/2}}
\newcommand{\RH}{\mathrm{RH}}

%% ─── Metadata ────────────────────────────────────────────────────────────────
\title{\textbf{Sturm--Liouville Eigenvalue Simplicity and a\\[4pt]
  Spectral Assumption Interface in Lean~4}\\[6pt]
  \large\textit{The Messina Nullification Framework}}

\author{
  Joseph Messina\\[4pt]
  \small Independent Researcher\\[2pt]
  \small \href{https://doi.org/10.5281/zenodo.18902325}{Zenodo: 10.5281/zenodo.18902325}
  \quad
  \href{https://github.com/Joeyxyxyz/messina-nullification-framework}{GitHub}
}

\date{March 2026}

%% ─────────────────────────────────────────────────────────────────────────────
\begin{document}
\maketitle

\begin{abstract}
We present two Lean~4/\Mathlib{} contributions.
The primary result is a formalization of Sturm--Liouville eigenvalue simplicity
for confining potentials: given $V(x) \to +\infty$ and a square-integrable
eigenfunction $\psi$, the Wronskian of any two eigenfunctions with the same
eigenvalue vanishes at infinity (by $L^2$ decay under cocompact convergence),
implying linear dependence and hence a one-dimensional eigenspace.
The argument is formalized using Mathlib's Picard--Lindel\"{o}f, $L^2$, and
cocompact filter infrastructure, and constitutes a standalone library
contribution independent of any number-theoretic application.
Second, building on the Lean~4 \Mathlib{} formalization of the Riemann zeta
function due to Loeffler and Stoll~\cite{loeffler2025}, we introduce a modular
spectral assumption interface under which the Riemann Hypothesis~(RH) and the
Simple Zeros Conjecture follow as conditional consequences of two explicitly
named hypotheses, with all classical analytic ingredients supplied by \Mathlib{}.
A constructive mechanism, the \emph{leveling trick}, shows that a spectral gap
at $\Delta = 1/2$ pins every nontrivial zero to the critical line via an
algebraic squeeze: the deviation $\re(\rho) - 1/2$ lies in a set with the
gap property, and is bounded in absolute value by $1/2$ from the critical
strip, so it must equal zero.
The single remaining open step --- connecting the operator's derived gap to the
shifted zeta-zero spectrum --- is isolated as a two-part named context field
\lean{HPContext} (object identification and operator identification), making the
Hilbert--P\'{o}lya correspondence explicit, modular, and independently
replaceable by future Mathlib developments.
All results are machine-checked under \texttt{leanprover/lean4:v4.24.0}
(Mathlib commit \texttt{f897ebcf}), zero-\lean{sorry}, with full code available
on Zenodo and GitHub.
The load-bearing algebraic content of every section has been independently
kernel-verified by AXLE (\emph{Axiom Lean Engine}); ten verification IDs
are listed in Section~\ref{sec:axle}.

\medskip\noindent
\textbf{2020 Mathematics Subject Classification:}
11M26 (Riemann zeta function),
34B24 (Sturm--Liouville theory),
03B35 (Mechanization of proofs).

\medskip\noindent
\textbf{Keywords:}
formal verification,
Lean~4,
Sturm--Liouville,
spectral gap,
Riemann zeta function,
Hilbert--P\'{o}lya correspondence,
Mathlib

\end{abstract}

%% ─────────────────────────────────────────────────────────────────────────────
\tableofcontents
\bigskip

%% ─────────────────────────────────────────────────────────────────────────────
\section{Introduction}
\label{sec:intro}

Recent advances in the \Lean{} theorem prover and the \Mathlib{} library have
enabled substantial portions of classical analysis to be formalized in a
machine-verified setting~\cite{mathlib2020}.
Formalization efforts have focused both on reproducing established results and on
developing reusable infrastructure for future analytic and spectral arguments.
Of particular relevance, Loeffler and Stoll recently formalized the Riemann zeta
function, the completed zeta function $\xi(s)$, and key analytic properties in
\Lean{}~\cite{loeffler2025}; the present work builds directly on that infrastructure.
In this paper we contribute two such developments.

\paragraph{Sturm--Liouville simplicity (unconditional contribution).}
We formalize a classical Sturm--Liouville simplicity result for one-dimensional
confining potentials.
Using \Mathlib{}'s analytic and functional-analytic infrastructure, we prove that
square-integrable eigenfunctions of confining Sturm--Liouville operators correspond
to simple eigenvalues.
The formalization proceeds through the standard Wronskian argument: constancy of
the Wronskian for solutions of the associated differential equation, asymptotic
$L^2$ decay of eigenfunctions at infinity, and the resulting vanishing of the
Wronskian implying linear dependence.
To our knowledge, this simplicity result is not currently available in
\Mathlib{} in a packaged, reusable form.

\begin{theorem}[Sturm--Liouville Simplicity, informal]
\label{thm:sl_informal}
Let $V \colon \R \to \R$ be a confining potential
($V$ continuous and $V(x) \to +\infty$ as $|x| \to \infty$).
Then every eigenvalue of the Sturm--Liouville operator $-d^2/dx^2 + V$
admitting a square-integrable eigenfunction is simple.
\end{theorem}

The Lean~4 statement is:
\begin{lstlisting}[language=Lean4]
theorem sturm_liouville_simple_eigenvalues (V : ℝ → ℝ)
    (hV : IsConfiningPotential V) (eigval : ℝ) :
    IsSimpleEigenvalue V eigval
\end{lstlisting}

\paragraph{Spectral dependency framework (conditional contribution).}
We introduce a modular spectral framework in \Lean{} for expressing conditional
implications related to the zeros of the Riemann zeta function.
The framework isolates a small set of explicitly named hypotheses --- including a
spectral gap condition and an identification principle relating formal spectral
objects to the shifted zeta-zero spectrum --- under which the Riemann Hypothesis
and simplicity of nontrivial zeros follow as formal consequences.
The purpose of this framework is \emph{not} to prove these hypotheses, but to
formalize the dependency structure underlying spectral approaches to the zeta zeros.
By expressing assumptions as explicit context fields, the formalization clarifies
precisely which analytic ingredients remain external to the Lean development.

The operator-algebraic component is built around an abstract projection-like
operator satisfying three structural properties: conservation ($\mathcal{N} + R = \mathrm{id}$),
idempotence ($\mathcal{N} \circ \mathcal{N} = \mathcal{N}$), and discreteness
(image $\subseteq \{0\} \cup [\Delta, \infty)$).
These properties enforce a spectral gap on the operator's range independent of
analytic number-theoretic arguments.
When combined with the identification hypothesis (Hilbert--P\'{o}lya correspondence,
stated as a named context field \lean{HPContext} in Section~\ref{sec:hpcontext}),
the gap condition yields RH and simplicity of nontrivial zeros as conditional consequences.

\begin{theorem}[Conditional RH and Simplicity, informal]
\label{thm:conditional_rh}
Under the Hilbert--P\'{o}lya identification hypothesis \lean{HPContext} and the
spectral gap property \lean{SpectralGapProperty(ShiftedZeroSpectrum, 1/2)},
all nontrivial zeros of $\zeta(s)$ lie on the critical line $\re(s) = 1/2$
and have order exactly~$1$.
\end{theorem}

\noindent
The spectral gap hypothesis should not be interpreted as strictly weaker than
the Riemann Hypothesis.
Rather, the framework isolates it as the single geometric constraint required
once the operator algebra is fixed; its connection to the actual zeta spectrum
requires the Hilbert--P\'{o}lya identification, which is a Millennium-class
open problem (Section~\ref{sec:obstruction}).

\paragraph{Claims we do not make.}
This paper does not claim a proof of the Riemann Hypothesis or the Simple Zeros
Conjecture.
It provides (i) an unconditional Sturm--Liouville library contribution and (ii)
a formal dependency analysis of spectral approaches to RH, with all assumptions
explicitly labeled and isolated.
The framework is designed as a reusable scaffold into which future Mathlib
analytic developments can be plugged to discharge remaining assumptions.

\paragraph{Reproducibility.}
The full environment is pinned and detailed in the Reproducibility
section (p.~\pageref{sec:reproducibility}).
All proofs compile under \texttt{leanprover/lean4:v4.24.0}
(Mathlib commit \texttt{f897ebcf}) with zero \lean{sorry}.

\paragraph{Contributions.}
\begin{enumerate}[label=\textbf{C\arabic*.}]
  \item \textbf{Sturm--Liouville simplicity (unconditional):}
        A packaged \Lean{} formalization of $\mathrm{IsConfiningPotential}
        \Rightarrow \mathrm{IsSimpleEigenvalue}$ via a Wronskian chain,
        independent of any RH-related application (Section~\ref{sec:proved}).
  \item \textbf{Formal dependency analysis (conditional):}
        A modular \Lean{} framework expressing $\RH \wedge \text{AllZerosSimple}$
        as consequences of an explicit, minimal assumption interface, with all
        classical ingredients supplied by \Mathlib{} (Section~\ref{sec:minimal}).
  \item \textbf{The leveling trick} (Section~\ref{sec:leveling}):
        a constructive proof that a spectral gap pins every nontrivial zero
        to the critical line via a two-step algebraic squeeze.
  \item \textbf{Assumption hygiene:}
        The Hilbert--P\'{o}lya identification is surfaced as a two-part named
        context field \lean{HPContext}, separating object identification from
        operator identification and making each assumption independently replaceable
        (Section~\ref{sec:hpcontext}).
  \item \textbf{Axiom reduction record:}
        A documented journey from 7 to 2 hypotheses characterizing precisely
        which classical results are absorbed by \Mathlib{} and which remain open
        (Section~\ref{sec:reduction}).
\end{enumerate}

\paragraph{Structure of the paper.}
Section~\ref{sec:operator} defines the nullification operator and its algebraic
properties.
Section~\ref{sec:setup} describes the formalization infrastructure.
Section~\ref{sec:proved} develops the Sturm--Liouville simplicity formalization in full, including definitions, the main theorem with Wronskian proof, and supporting Lean infrastructure.
Section~\ref{sec:leveling} develops the leveling trick.
Section~\ref{sec:minimal} states the main conditional reduction theorem.
Section~\ref{sec:reduction} documents the hypothesis reduction journey.
Section~\ref{sec:obstruction} characterizes the irreducible obstruction and
introduces \lean{HPContext}.
Section~\ref{sec:lessons} reflects on mathematical lessons from formalization.
Section~\ref{sec:conclusion} concludes.
Appendix~\ref{app:motivation} contains motivational comparisons with
non-standard analysis and a concrete numerical example of non-distributivity.
Appendix~\ref{app:experimental} describes the experimental companion artifact
on IBM quantum hardware (non-deductive motivation only, not required for any
formal result).

%% ─────────────────────────────────────────────────────────────────────────────
\section{The Nullification Operator and Its Algebraic Properties}
\label{sec:operator}

The nullification operator arises as an abstract projection-like map
enforcing a discrete spectral partition: values below a threshold $\Delta$
are sent to zero (``confined'' or ``nullified''), while values at or above
threshold pass through unchanged.
Its three algebraic properties --- conservation, idempotence, and discreteness
--- are sufficient to enforce a gap on the operator's range, independent of
any analytic number-theoretic argument.

\begin{definition}[Nullification operator]
\label{def:nullify}
For threshold $\Delta > 0$ and value $x \in \R$, define:
\[
  \N(\Delta, x) \;:=\;
  \begin{cases}
    x & \text{if } |x| \geq \Delta, \\
    0 & \text{if } |x| < \Delta.
  \end{cases}
\]
\end{definition}

Every value is assigned to one of three \emph{classes}:
\[
  [0] \;:\; x = 0,
  \qquad
  [1] \;:\; |x| \geq \Delta,
  \qquad
  [0.0e] \;:\; 0 < |x| < \Delta.
\]
The class label $[0.0e]$ — read informally as a small value that
``rounds to zero in observation but retains a nonzero magnitude'' — is
chosen to suggest both an \emph{exponentially-suppressed} value and an
\emph{epsilon-small} one; the formal content is the inequality
$0 < |x| < \Delta$ and is independent of either reading.
The class $[0.0e]$ — ``nullified'' or ``confined'' — represents values
below the propagation threshold that are not absent but preserved: the companion
operator $R(\Delta, x) := x - \N(\Delta, x)$ routes sub-threshold content
to a \emph{remainder pool} rather than discarding it.
The conservation law $\N(\Delta, x) + R(\Delta, x) = x$ holds for all $x$.

The operator satisfies strict conservation:

\begin{proposition}[Conservation]
\label{prop:conservation}
For all $\Delta > 0$ and $x \in \R$:
$\N(\Delta, x) + R(\Delta, x) = x$,
where $R(\Delta, x) := x - \N(\Delta, x)$.
\end{proposition}

Four key algebraic properties are proved in the Batch series:

\begin{enumerate}[label=\textbf{P\arabic*.}]
  \item \textbf{Idempotence:} $\N(\Delta, \N(\Delta, x)) = \N(\Delta, x)$.
  \item \textbf{Monotonicity:} $\N(\Delta, \cdot)$ is monotone above threshold.
  \item \textbf{Non-distributivity:}
        $\N(\Delta, x + y) \neq \N(\Delta, x) + \N(\Delta, y)$ in general.
        This is not a defect: it captures threshold-crossing effects.
  \item \textbf{Spectral gap:} For a set $S \subseteq \R$, after applying $\N$,
        the interval $(0, \Delta)$ is empty in the image.
\end{enumerate}

The spectral gap property is the central structural concept:

\begin{definition}[Spectral gap property]
\label{def:sgp}
A set $S \subseteq \R$ has the \emph{spectral gap property at $\Delta$}
if every $x \in S$ satisfies $x = 0$ or $|x| \geq \Delta$:
\[
  \mathrm{SpectralGapProperty}(S, \Delta)
  \;:\Longleftrightarrow\;
  \forall\, x \in S,\; x = 0 \;\vee\; |x| \geq \Delta.
\]
\end{definition}

In Lean~4 this is stated as:

\begin{lstlisting}[language=Lean4]
def SpectralGapProperty (S : Set ℝ) (Δ : ℝ) : Prop :=
  ∀ x, x ∈ S → x = 0 ∨ |x| ≥ Δ
\end{lstlisting}

In subsequent statements, the set $S$ is instantiated as
\lean{ShiftedZeroSpectrum} (Section~\ref{sec:setup}); the threshold
$\Delta$ is instantiated as $1/2$ throughout the leveling-trick development.

%% ─────────────────────────────────────────────────────────────────────────────
\subsection{Spectral Gap Enforcement}
\label{sec:gap_enforcement}

The three algebraic properties — conservation, idempotence, and discreteness
— are not merely structural observations; they jointly enforce a spectral gap
on the operator's range.

\begin{proposition}[\lean{master\_enforcement}]
\label{prop:enforcement}
For any $\Delta > 0$, the image of $\N_\Delta$ satisfies
$\mathrm{SpectralGapProperty}(\mathrm{Set.range}(\N_\Delta), \Delta)$.
That is, every element of the range equals zero or has absolute value at
least $\Delta$.
\end{proposition}

\begin{proof}
By the definition of $\N(\Delta, x)$: if $|x| \geq \Delta$ then
$\N(\Delta, x) = x$ with $|x| \geq \Delta$; if $|x| < \Delta$ then
$\N(\Delta, x) = 0$.
No element of the range lies in the open interval $(0, \Delta)$.
In Lean~4 (Batch~14, \lean{master\_enforcement}):
\begin{lstlisting}[language=Lean4]
theorem master_enforcement (Δ : ℝ) (hΔ : 0 < Δ) :
    SpectralGapProperty (Set.range (N Δ)) Δ := by
  intro x hx
  obtain ⟨y, rfl⟩ := hx
  by_cases h : |y| ≥ Δ
  · -- |y| ≥ Δ: N Δ y = y, so |x| = |y| ≥ Δ; right branch
    right
    simpa [N, h] using h
  · -- |y| < Δ: N Δ y = 0; left branch
    left
    simp [N, h]
\end{lstlisting}
The two branches correspond to the two cases of the definition of $\N$:
when the input is at or above threshold the value passes through and
satisfies the gap bound; when it is below threshold the value is sent
to zero. The pattern \lean{by\_cases} (rather than \lean{split\_ifs})
makes the case analysis explicit and the branch order clear.
\end{proof}

\begin{remark}
This gap is derived purely from operator algebra: no analytic number theory
is used.
The step that requires the Hilbert--P\'{o}lya identification is the transfer
of this gap from $\mathrm{Set.range}(\N_\Delta)$ to the actual
shifted zeta-zero spectrum $\mathrm{ShiftedZeroSpectrum}$
(Section~\ref{sec:hpcontext}).
\end{remark}

%% ─────────────────────────────────────────────────────────────────────────────
\begin{remark}[On comparisons with non-standard analysis]
\label{rem:nsa-pointer}
We emphasize that $\N(\Delta, \cdot)$ is not a ring homomorphism and is
intentionally non-distributive; conservation holds via the companion
remainder map $R(\Delta, x) := x - \N(\Delta, x)$ as
$\N(\Delta, x) + R(\Delta, x) = x$.
A more extended comparison with Robinson's standard part function
$\mathrm{st}()$ --- including a concrete numerical example of
non-distributivity --- is deferred to
Appendix~\ref{app:motivation}, since these comparisons are motivational
rather than formal contributions of this paper.
\end{remark}

%% ─────────────────────────────────────────────────────────────────────────────
\section{Formalization Setup and Key Definitions}
\label{sec:setup}

\paragraph{Infrastructure.}
The formalization uses \Lean{} version \texttt{leanprover/lean4:v4.24.0}
with Mathlib commit \texttt{f897ebcf}.
All analytic properties of the Riemann zeta function --- analyticity,
the functional equation, the completed zeta function $\xi(s)$,
and the Schwarz reflection principle --- are imported directly from
Loeffler--Stoll~\cite{loeffler2025} via Mathlib.
The \Aristotle{} automated theorem prover~\cite{aristotle2025}
was used to close proof gaps in the formalization.

\paragraph{Key definitions.}

\begin{lstlisting}[language=Lean4]
-- Nontrivial zeros: zeros in the critical strip
def is_nontrivial_zero (rho : Complex) : Prop :=
  riemannZeta rho = 0 and 0 < rho.re and rho.re < 1

-- Order of vanishing via Mathlib infrastructure
noncomputable def zeta_order (rho : Complex) : Nat :=
  ENat.toNat (AnalyticAt.order riemannZeta rho)

-- Shifted zero spectrum: deviations from Re = 1/2
def ShiftedZeroSpectrum : Set Real :=
  {x | exists rho in RiemannZeros, x = rho.re - 1/2}

-- The Riemann Hypothesis
def RiemannHypothesis : Prop :=
  forall rho : Complex, is_nontrivial_zero rho -> rho.re = 1/2

-- All zeros simple
def AllZerosSimple : Prop :=
  forall rho : Complex, is_nontrivial_zero rho -> zeta_order rho = 1
\end{lstlisting}

\paragraph{Scale.}
The formalization spans 82~Lean~4 files organized in three volumes:
\begin{enumerate}
  \item \textbf{Core operator library} (Batch 1--14, MdBatch 1--7):
        algebraic properties of $\N(\Delta,x)$ in full generality.
  \item \textbf{RH and number theory} (Outputs 0--86):
        the reduction machinery, leveling trick, and minimal axiom surface.
  \item \textbf{Extended applications} (Outputs 87--96):
        Sturm--Liouville bridge to the mass gap problem and related spectral systems.
\end{enumerate}

%% ─────────────────────────────────────────────────────────────────────────────
\section{Sturm--Liouville Simplicity Formalization}
\label{sec:proved}

\subsection{Main SL Simplicity Theorem}
\label{sec:sl_main}

The Sturm--Liouville simplicity result is the primary unconditional contribution
of this paper.
We formalize the classical Wronskian argument for confining one-dimensional
Sturm--Liouville operators.
The key definitional components are:

\begin{lstlisting}[language=Lean4]
-- A potential V is confining: continuous and V(x) → +∞ as |x| → ∞
def IsConfiningPotential (V : ℝ → ℝ) : Prop :=
  Continuous V ∧ Filter.Tendsto V (Filter.cocompact ℝ) Filter.atTop

-- u is an eigenfunction of -d²/dx² + V with eigenvalue eigval,
-- nonzero, and square-integrable
def IsSturmLiouvilleEigenfunction
    (V : ℝ → ℝ) (eigval : ℝ) (u : ℝ → ℝ) : Prop :=
  Differentiable ℝ u ∧
  Differentiable ℝ (deriv u) ∧
  (∀ x, -(deriv (deriv u) x) + V x * u x = eigval * u x) ∧
  u ≠ 0 ∧
  MeasureTheory.MemLp u 2 MeasureTheory.volume

-- The Wronskian of two functions u, v at a point x
def wronskian (u v : ℝ → ℝ) (x : ℝ) : ℝ :=
  u x * deriv v x - v x * deriv u x

-- An eigenvalue is simple if every two eigenfunctions
-- for it are scalar multiples of each other
def IsSimpleEigenvalue (V : ℝ → ℝ) (eigval : ℝ) : Prop :=
  ∀ u v : ℝ → ℝ,
    IsSturmLiouvilleEigenfunction V eigval u →
    IsSturmLiouvilleEigenfunction V eigval v →
    ∃ c : ℝ, ∀ x, v x = c * u x
\end{lstlisting}

\begin{theorem}[Sturm--Liouville Simplicity, \lean{sturm\_liouville\_simple\_eigenvalues}]
\label{thm:sl_simplicity}
Let $V \colon \R \to \R$ be a confining potential.
Then every eigenvalue of $-d^2/dx^2 + V$ admitting a square-integrable
eigenfunction is simple: any two such eigenfunctions for the same
eigenvalue are scalar multiples of each other.
\end{theorem}

\begin{proof}
Fix an eigenvalue $\mathit{eigval}$ and two eigenfunctions $u, v$ for it.
The Wronskian $W[u,v](x) = u(x)v'(x) - v(x)u'(x)$ of two solutions to
the same eigenvalue equation is constant (Abel's identity, packaged as
\lean{wronskian\_constant\_same\_eigenvalue}):
\[
  \frac{d}{dx} W[u,v] = 0 \quad \Longrightarrow \quad W \equiv C \in \R.
\]
If both $u, v \in L^2(\R)$ and $V$ is confining, then
$|u(x)|, |v(x)| \to 0$ as $|x| \to \infty$, and likewise for $u', v'$
via the second-order ODE
$u'' = (V(x) - \mathit{eigval})\, u(x)$.
Hence $W[u,v](x) \to 0$ as $|x| \to \infty$, so $C = 0$, so $W \equiv 0$
(\lean{wronskian\_zero\_of\_l2\_final}).
$W \equiv 0$ together with $u \neq 0$ then forces $v$ to be a scalar
multiple of $u$ (\lean{linearly\_dependent\_of\_wronskian\_zero}).

In Lean~4 (pinned to \texttt{f897ebcf}), the capstone reads:
\begin{lstlisting}[language=Lean4]
theorem sturm_liouville_simple_eigenvalues (V : ℝ → ℝ)
    (hV : IsConfiningPotential V) (eigval : ℝ) :
    IsSimpleEigenvalue V eigval := by
  intro u v hu hv
  have hW : ∀ x, wronskian u v x = 0 :=
    wronskian_zero_of_l2_final V eigval hV u v hu hv
  exact linearly_dependent_of_wronskian_zero
    V eigval hV.1 u v hu hv hW
\end{lstlisting}
The two intermediate lemmas \lean{wronskian\_zero\_of\_l2\_final} and
\lean{linearly\_dependent\_of\_wronskian\_zero} are themselves proved
by combining \lean{wronskian\_constant\_same\_eigenvalue} (Abel's
identity) with the asymptotic decay
\lean{tendsto\_zero\_and\_deriv\_zero\_of\_eigenfunction\_final} for
$L^2$ eigenfunctions of confining potentials.
\end{proof}

\subsection{Supporting Theorems}
\label{sec:supporting_thms}

The following results, gathered in \lean{MinimalAxiomReduction.lean},
are established with zero unresolved \lean{sorry} placeholders, zero
\lean{exact?} fallbacks, and zero \lean{admit} calls.
All proofs compile against Mathlib.
The Sturm--Liouville capstone of Section~\ref{sec:sl_main} also compiles
zero-\lean{sorry} and zero-\lean{admit}; it contains a small number of
\lean{exact?} calls produced during \Aristotle{}-assisted gap closure,
each of which the Lean kernel type-checks to a concrete proof term at
elaboration time.

\begin{theorem}[\lean{shifted\_bounded}]
\label{thm:shifted_bounded}
For every non-trivial zero $\rho$ of $\zeta(s)$:
$|\re(\rho) - 1/2| < 1/2$.
\end{theorem}

\begin{proof}
Immediate from $0 < \re(\rho) < 1$ (critical strip), algebraic manipulation.
In Lean:
\begin{lstlisting}[language=Lean4]
theorem shifted_bounded :
    forall x in ShiftedZeroSpectrum, |x| < 1/2 := by
  intro x hx
  obtain <rho, hrho, rfl> := hx
  rw [abs_lt]
  exact <by linarith [hrho.2.1], by linarith [hrho.2.2]>
\end{lstlisting}
\end{proof}

\begin{theorem}[\lean{tau\_half\_real}]
\label{thm:tau_half}
If $S \subseteq \R$ satisfies the spectral gap property at $\Delta$
and every element of $S$ satisfies $|x| < \Delta$,
then $S \subseteq \{0\}$.
\end{theorem}

\begin{proof}
For $x \in S$: the gap property gives $x = 0$ or $|x| \geq \Delta$;
the bound $|x| < \Delta$ excludes the second case. Hence $x = 0$.
\end{proof}

\begin{theorem}[\lean{zeta\_order\_ge\_one}]
\label{thm:order_ge_one}
Every non-trivial zero of $\zeta(s)$ has order at least~$1$.
\end{theorem}

\begin{proof}
A zero of an analytic function has order $\geq 1$ by definition.
Proved via \lean{AnalyticAt.order\_eq\_zero\_iff} from Mathlib.
\end{proof}

\begin{theorem}[\lean{order\_one\_iff\_deriv}]
\label{thm:order_iff}
For a non-trivial zero $\rho$:
$\zetaord_\rho(\zeta) = 1 \iff \zeta'(\rho) \neq 0$.
\end{theorem}

\begin{proof}
The order of vanishing of an analytic function at a zero equals~$1$
precisely when the first derivative is nonzero at that point.
Proved via \lean{AnalyticAt.order\_eq\_one\_iff} applied to
\lean{riemannZeta\_analyticAt} from Mathlib:
\begin{lstlisting}[language=Lean4]
theorem order_one_iff_deriv
    (hrho : is_nontrivial_zero rho) :
    zeta_order rho = 1 <-> deriv riemannZeta rho != 0 :=
  AnalyticAt.order_eq_one_iff (riemannZeta_analyticAt rho)
\end{lstlisting}
\end{proof}

\begin{theorem}[\lean{xi\_even}]
\label{thm:xi_even}
The completed zeta function satisfies
$\xi(1/2 + w) = \xi(1/2 - w)$ for all $w \in \C$.
\end{theorem}

\begin{proof}
Direct from \lean{completedRiemannZeta\_one\_sub} in Mathlib via \lean{ring\_nf}.
\end{proof}

\begin{theorem}[\lean{GUE\_level\_repulsion}]
\label{thm:gue}
The GUE pair correlation function satisfies $R_2(0) = 0$.
\end{theorem}

\begin{proof}
By definition of $R_2(\alpha) = 1 - (\sin(\pi\alpha)/(\pi\alpha))^2$
for $\alpha \neq 0$ and $R_2(0) = 0$: \lean{simp [GUE\_pair\_correlation]}.
\end{proof}

\begin{theorem}[\lean{multiplicity\_squeeze}]
\label{thm:squeeze}
If every non-trivial zero has order $\geq 1$ and order $\leq 1$,
then every non-trivial zero has order exactly~$1$.
\end{theorem}

\begin{proof}
Immediate from \lean{Nat.le\_antisymm}.
\end{proof}

%% ─────────────────────────────────────────────────────────────────────────────
\section{The Leveling Trick}
\label{sec:leveling}

The most structurally novel contribution of this formalization is the
\emph{leveling trick} --- a constructive mechanism that pins every
non-trivial zero to the critical line $\re(s) = 1/2$.

\subsection{Algebraic motivation}

The algebraic content of the leveling trick is the observation that any
real number $|x| < 1/2$ in a set with the spectral gap property at
$\Delta = 1/2$ must be exactly zero: the gap forces $x = 0$ or $|x|
\geq 1/2$, and the bound $|x| < 1/2$ rules out the second alternative.
Applied to the deviation $\re(\rho) - 1/2$ of a nontrivial zero, this
forces $\re(\rho) = 1/2$.

\begin{definition}[Signed deviation and its conjugate]
For a point $\rho \in \C$, define:
\begin{align*}
  \mathrm{Dev}(\rho)      &:= \re(\rho) - 1/2, \\
  \mathrm{ConjDev}(\rho)  &:= -(\re(\rho) - 1/2) = 1/2 - \re(\rho).
\end{align*}
\end{definition}

These are antisymmetric: $\mathrm{ConjDev}(\rho) = -\mathrm{Dev}(\rho)$,
so $\mathrm{Dev}(\rho) + \mathrm{ConjDev}(\rho) = 0$ identically.
Note that $\mathrm{Dev}(\rho) \in \mathrm{ShiftedZeroSpectrum}$ for every
nontrivial zero $\rho$.

\subsection{The balance condition}

The functional equation $\zeta(s) = \zeta(1-s)$ enforces a balance
between $\rho$ and its conjugate zero $1 - \rho$:
\[
  \mathrm{Dev}(1 - \rho)
  = \re(1 - \rho) - \tfrac{1}{2}
  = 1 - \re(\rho) - \tfrac{1}{2}
  = -\mathrm{Dev}(\rho).
\]
So $\mathrm{Dev}(1-\rho) = -\mathrm{Dev}(\rho)$: functional-equation partners
have opposite signed deviations.
A single zero self-paired under the functional equation (i.e., $1-\rho = \rho$)
must therefore satisfy $\mathrm{Dev}(\rho) = -\mathrm{Dev}(\rho)$,
which forces $\mathrm{Dev}(\rho) = 0$, i.e., $\re(\rho) = 1/2$.

\subsection{The leveling theorem}

\begin{theorem}[Leveling theorem, \lean{leveling\_theorem}]
\label{thm:leveling}
If $\mathrm{ShiftedZeroSpectrum}$ has the spectral gap property at $\Delta = 1/2$,
then for every nontrivial zero $\rho$: $\mathrm{Dev}(\rho) = 0$,
i.e., $\re(\rho) = 1/2$.
\end{theorem}

\begin{proof}
Since $\rho$ is a nontrivial zero, $\mathrm{Dev}(\rho) = \re(\rho) - 1/2
\in \mathrm{ShiftedZeroSpectrum}$.
By the gap property: either $\mathrm{Dev}(\rho) = 0$ or
$|\mathrm{Dev}(\rho)| \geq 1/2$.
By Theorem~\ref{thm:shifted_bounded}, $|\re(\rho) - 1/2| < 1/2$,
so the second case is impossible.
Hence $\mathrm{Dev}(\rho) = 0$ and $\re(\rho) = 1/2$.

In Lean~4:
\begin{lstlisting}[language=Lean4]
theorem leveling_theorem
    (h_gap : SpectralGapProperty ShiftedZeroSpectrum (1/2))
    (rho : ℂ) (hrho : is_nontrivial_zero rho) :
    Dev rho = 0 := by
  have hmem : Dev rho ∈ ShiftedZeroSpectrum := ⟨rho, hrho, rfl⟩
  rcases h_gap (Dev rho) hmem with h0 | hge
  · exact h0
  · exfalso; linarith [shifted_bounded hrho, hge]
\end{lstlisting}
\end{proof}

\begin{corollary}[\lean{spectral\_gap\_implies\_rh}]
\label{cor:gap_rh}
$\mathrm{SpectralGapProperty}(\mathrm{ShiftedZeroSpectrum}, 1/2) \implies \RH$.
\end{corollary}

\begin{proof}
By Theorem~\ref{thm:leveling}: $\mathrm{Dev}(\rho) = 0$ for every nontrivial
zero, hence $\re(\rho) = 1/2$.
\end{proof}

\begin{remark}
The leveling theorem is a direct algebraic consequence of two facts:
the gap property forces values into $\{0\} \cup \{x : |x| \geq 1/2\}$,
and the bounded strip (Theorem~\ref{thm:shifted_bounded}) forces
$|\mathrm{Dev}(\rho)| < 1/2$.
The intersection is the single point $\{0\}$.
This is the structural content of ``leveling'': every individual nontrivial
zero is pinned to the critical line by the same two-step squeeze, in
contrast to asymptotic results that bound averages but do not pin individual
zeros.
\end{remark}

%% ─────────────────────────────────────────────────────────────────────────────
\section{The Minimal Axiom Reduction}
\label{sec:minimal}

\begin{definition}[Axiom surface]
\label{def:axioms}
The two named hypotheses forming the minimal assumption interface are:
\begin{align*}
  \mathbf{SpectralGap} &\;:\;
    \mathrm{SpectralGapProperty}(\mathrm{ShiftedZeroSpectrum},\; 1/2), \\
  \mathbf{SimpleBound} &\;:\;
    \forall\, \rho\;\text{non-trivial},\;
    \zetaord_\rho(\zeta) \leq 1.
\end{align*}
\end{definition}

\begin{theorem}[Main reduction]
\label{thm:main}
$\mathbf{SpectralGap} \wedge \mathbf{SimpleBound}
\;\implies\;
\RH \;\wedge\; \mathrm{AllZerosSimple}$.
\end{theorem}

\begin{proof}
\textbf{RH.} By Corollary~\ref{cor:gap_rh} applied to \textbf{SpectralGap}.

\textbf{AllZerosSimple.} By Theorem~\ref{thm:order_ge_one},
every non-trivial zero has order $\geq 1$.
By \textbf{SimpleBound}, every non-trivial zero has order $\leq 1$.
By Theorem~\ref{thm:squeeze}, every non-trivial zero has order $= 1$.

In Lean:
\begin{lstlisting}[language=Lean4]
-- The fully formalized capstone takes three hypotheses:
-- gap, bridge (discharged by the leveling theorem of §5), and bound.
theorem resolution_from_three_axioms
    (h_gap : SpectralGapProperty ShiftedZeroSpectrum (1/2))
    (h_bridge : ShiftedZeroSpectrum ⊆ {0} → RiemannHypothesis)
    (h_simple : MultiplicityBoundOne) :
    RiemannHypothesis ∧ AllZerosSimple := by
  refine ⟨?_, ?_⟩
  · exact h_bridge (tau_half_real _ shifted_bounded h_gap)
  · exact bound_implies_simple h_simple

-- Specializing the bridge to spectral_gap_implies_rh (the §5 leveling
-- corollary) gives the two-axiom form stated in Theorem 6.2:
theorem resolution_from_two_axioms
    (h_gap : SpectralGapProperty ShiftedZeroSpectrum (1/2))
    (h_simple : MultiplicityBoundOne) :
    RiemannHypothesis ∧ AllZerosSimple :=
  resolution_from_three_axioms h_gap
    (fun _ => spectral_gap_implies_rh h_gap)
    h_simple
\end{lstlisting}
The bridge premise $\mathrm{ShiftedZeroSpectrum} \subseteq \{0\} \to \mathrm{RiemannHypothesis}$
is itself an immediate consequence of the leveling theorem
(Section~\ref{sec:leveling}), so the two-axiom statement above is
equivalent to the three-axiom packaging in
\lean{MinimalAxiomReduction.lean}.
\end{proof}

The two axioms are also sufficient for each conjecture independently:

\begin{corollary}
\label{cor:independent}
\begin{enumerate}[label=(\alph*)]
  \item $\mathbf{SpectralGap} \implies \RH$
        (without simplicity, \lean{resolution\_rh\_only}).
  \item $\mathbf{SimpleBound} \implies \mathrm{AllZerosSimple}$
        (without RH, \lean{resolution\_simple\_only}).
\end{enumerate}
\end{corollary}

Four independent resolution paths are formalized:
\begin{center}
\begin{tabular}{lll}
  \toprule
  Path & Hypotheses & Conclusion \\
  \midrule
  A (Derivative) & SpectralGap + DerivativeNonvanishing & RH $\wedge$ Simple \\
  B (Multiplicity) & SpectralGap + MultBoundOne & RH $\wedge$ Simple \\
  C (Contrapositive) & SpectralGap + MultipleZeroKillsGap & RH $\wedge$ Simple \\
  D (GUE) & SpectralGap + GUE\_Exclusion & RH $\wedge$ Simple \\
  \bottomrule
\end{tabular}
\end{center}
Paths A and B are equivalent (Theorem~\ref{thm:order_iff}).
Paths B, C, and D are machine-verified.

\begin{remark}[On the strength of \textbf{SpectralGap}]
\label{rem:sg_strength}
The hypothesis \textbf{SpectralGap} is not claimed to be strictly weaker
than RH.
If \lean{ShiftedZeroSpectrum} is taken as the literal set of deviations
$\{\re(\rho) - 1/2 \mid \zeta(\rho) = 0,\; 0 < \re(\rho) < 1\}$,
then \textbf{SpectralGap} is logically equivalent to RH given the
bounded strip (Theorem~\ref{thm:shifted_bounded}).
The framework's contribution is not a weakening of RH but an
\emph{isolation}: the operator algebra of $\N(\Delta, x)$ (Batch~14,
\lean{master\_enforcement}) derives a gap property on
$\mathrm{Set.range}(\N_\Delta)$ from three algebraic properties
(conservation, idempotence, discreteness).
The single remaining open step is the identification
$\mathrm{ShiftedZeroSpectrum} = f(\mathrm{Set.range}(\N_\Delta))$
for an appropriate map $f$.
This identification --- the Hilbert--P\'{o}lya correspondence --- is
stated as an explicit named context field \lean{hp\_correspondence} in
Section~\ref{sec:obstruction}, separating all algebraic content from the
single open conjecture.
\end{remark}

%% ─────────────────────────────────────────────────────────────────────────────
\section{The Axiom Reduction Journey}
\label{sec:reduction}

A key lesson of this formalization is how much can be eliminated from the
axiom surface through careful interaction with \Mathlib{}.
The reduction proceeded in four stages:

\begin{center}
\begin{tabular}{llll}
  \toprule
  Output & Name & Axioms & Result \\
  \midrule
  40  & Capstone v1    & 7 & $\RH \wedge \text{Simple}$ \\
  55  & Weinstein form & 5 & $\RH \wedge \text{Simple}$ \\
  68  & Grand unified  & 5 & $\RH \wedge \text{Simple}$ \\
  71  & \textbf{Minimal}      & \textbf{2} & $\RH \wedge \text{Simple}$ \\
  71b & RH only        & 1 & $\RH$ \\
  71c & Simple only    & 1 & $\text{Simple}$ \\
  \bottomrule
\end{tabular}
\end{center}

Axioms eliminated in the reduction from 7 to 2:
\begin{enumerate}
  \item \textbf{Analyticity} of $\zeta$ --- available in Mathlib via
        \lean{riemannZeta\_analyticAt}.
  \item \textbf{Functional equation} $\zeta(s) = \zeta(1-s)$ --- Mathlib,
        \lean{completedRiemannZeta\_one\_sub}.
  \item \textbf{$\xi$ symmetry} $\xi(1/2+w) = \xi(1/2-w)$ --- proved as
        Theorem~\ref{thm:xi_even}, zero additional hypotheses.
  \item \textbf{Order $\geq 1$} at non-trivial zeros --- proved as
        Theorem~\ref{thm:order_ge_one} via \lean{AnalyticAt.order}.
  \item \textbf{Speiser's theorem} formalized in output~45 and imported.
\end{enumerate}

Each eliminated axiom corresponds to a theorem provable from Mathlib
infrastructure without additional number-theoretic input.
The two remaining axioms are of a different character, as we now explain.

%% ─────────────────────────────────────────────────────────────────────────────
\section{The Irreducible Obstruction}
\label{sec:obstruction}

\begin{theorem}[Irreducibility of SpectralGap]
\label{thm:sg_irreducible}
\textbf{SpectralGap} cannot be proved from Mathlib without solving
a Millennium-class problem.
\end{theorem}

\begin{proof}[Informal]
Making SpectralGap a theorem requires proving that
$\mathrm{ShiftedZeroSpectrum}$ has no elements in $(0, 1/2)$.
This is equivalent to proving that all non-trivial zeros satisfy $\re(\rho) = 1/2$,
i.e., the Riemann Hypothesis itself.
Alternatively, it requires proving that zeta zeros are eigenvalues of a
self-adjoint operator, which is the content of the Hilbert--P\'olya conjecture,
itself a Millennium-class problem.
\end{proof}

\begin{theorem}[Irreducibility of SimpleBound]
\label{thm:sb_irreducible}
\textbf{SimpleBound} is logically independent of $\RH$ and cannot be
eliminated by any known classical technique.
\end{theorem}

\begin{proof}[Informal]
We verify each known approach fails:
\begin{enumerate}
  \item \textbf{Laguerre--P\'olya class:}
        LP class membership permits double zeros by definition.
  \item \textbf{Rolle's theorem:}
        The zero count bound is exactly tight (giving $n-1$ zeros for
        degree~$n$), so Rolle provides no margin.
  \item \textbf{Speiser's theorem:}
        Speiser governs $\re(s) < 1/2$, not $\re(s) = 1/2$.
  \item \textbf{de Bruijn--Newman dynamics:}
        The flow degenerates at $\Lambda = 0$, providing no information
        about zero order.
  \item \textbf{Goldston--Gonek--Montgomery (2025):}
        Pair correlation implies asymptotically 100\% simple zeros,
        but ``asymptotically all'' $\neq$ ``all.''
\end{enumerate}
The conjecture remains open.
\end{proof}

\begin{remark}
The density-zero obstruction is the precise gap.
The set of non-simple zeros could have density zero among all zeros
(consistent with Goldston et al.) yet still be infinite.
The conjecture \textbf{SimpleBound} is precisely the statement that this
density-zero set is empty.
\end{remark}

%% ─────────────────────────────────────────────────────────────────────────────
\subsection{The Identification Context: \lean{HPContext}}
\label{sec:hpcontext}

Spectral approaches to the Riemann Hypothesis, such as the
Hilbert--P\'{o}lya program, posit the existence of a self-adjoint operator
whose spectrum corresponds to the imaginary parts of the nontrivial zeros of
$\zeta(s)$~\cite{montgomery1973,goldston2001}.
If such an operator exists, its eigenvalues lie on the real axis by
self-adjointness, placing the associated zeros on $\re(s) = 1/2$ --- exactly
the content of RH.
The single gap in the present framework that cannot be closed by operator
algebra alone is the identification of $\N$'s derived spectral range with the
actual zeta-zero shift spectrum.
To make this identification explicit and modular --- separating \emph{object
identification} from \emph{operator identification} --- the Lean formalization
introduces:

\begin{lstlisting}[language=Lean4]
-- The load-bearing bridge, surfaced as a named context.
-- Two separate identification steps prevent conflation of distinct assumptions:
structure HPContext (Delta : Real) where
  -- (1) Object identification: what spectrum corresponds to zeta zeros
  zetaShift : Complex -> Real         -- Re(rho) - 1/2
  zeroSet   : Set Complex             -- nontrivial zeta zeros
  shiftedSpec : Set Real
  shiftedSpec_def : shiftedSpec = Set.image zetaShift zeroSet

  -- (2) Operator identification: N acts on the right object  
  N_targets : Set.image zetaShift zeroSet ⊆ Set.range (N Delta)
  -- Consequence: SpectralGapProperty (Set.range (N Delta)) Delta
  -- + N_targets gives SpectralGapProperty shiftedSpec Delta
  -- The operator-algebra gap (Batch14, master_enforcement) then transfers.
\end{lstlisting}

This structure makes the Hilbert--P\'{o}lya correspondence a
\emph{first-class hypothesis with explicit components}, not an implicit
background assumption.
The spectral gap property should not be interpreted as strictly weaker than
the Riemann Hypothesis: under the literal definition of
\lean{ShiftedZeroSpectrum}, the two are equivalent by the bounded strip
(Theorem~\ref{thm:shifted_bounded}).
What the framework contributes is \emph{isolation}: the algebraic content
of $\N$ (conservation, idempotence, discreteness) is proved independently,
and the gap's connection to the actual zeta spectrum is precisely the
single open conjecture surfaced by \lean{HPContext}.
Future work that constructs a self-adjoint operator whose spectrum
matches zeta zeros (the classical Hilbert--P\'{o}lya program) would
instantiate \lean{HPContext} with proved theorems, at which point all
\lean{sorry}-free theorems in this framework discharge automatically.

%% ─────────────────────────────────────────────────────────────────────────────
\section{Mathematical Lessons from Formalization}
\label{sec:lessons}

\paragraph{Lesson 1: Mathlib absorbs more than expected.}
Of the seven initial axioms, five were eliminable through careful
interaction with Mathlib's existing zeta function infrastructure.
The analyticity, functional equation, and xi symmetry --- each treated as
independent axioms in earlier versions --- all followed from theorems
already formalized by Loeffler--Stoll~\cite{loeffler2025}.
This underscores the value of building on established Mathlib infrastructure
rather than re-axiomatizing standard results.

\paragraph{Lesson 2: The axiom surface is the proof.}
The journey from 7 to 2 axioms is not merely a technical optimization:
it reveals the precise logical structure of the problem.
Two named hypotheses remain because they correspond to two genuinely
open problems with no known proof pathway.
The formalization makes this structure \emph{explicit and checkable}:
one can inspect exactly which theorems depend on which assumptions.

\paragraph{Lesson 3: Automated provers as completion engines.}
The \Aristotle{} theorem prover~\cite{aristotle2025} was used extensively
for gap closure.
Correctness is ensured by the Lean kernel, which type-checks every
proof term regardless of how it was found.
The kernel is the arbiter of truth; \Aristotle{} acts as a search
heuristic over candidate proof terms, not as an independent semantic
authority.
In practice, careful prompt engineering and explicit definitional
control were nonetheless required to keep \Aristotle{} aimed at the
intended theorem statements: in novel frameworks that do not match
conventional patterns, human conceptual direction must remain primary,
with automation handling technical gap-filling under kernel
verification.

\paragraph{Lesson 4: Algebraic motivation guides formalization.}
The leveling theorem arose from a simple algebraic observation: a value
in a set with the spectral gap property at $\Delta = 1/2$ must satisfy
$x = 0$ or $|x| \geq 1/2$, while the bounded strip forces
$|\mathrm{Dev}(\rho)| < 1/2$.
The intersection is the single point $\{0\}$.
Recognizing this two-step squeeze --- gap plus bound --- before writing
the Lean proof determined the structure of \lean{leveling\_theorem}:
case-split on the gap disjunction, then close the second case with a
single \lean{linarith} call from the strip bound.
The same algebraic pattern appears throughout the corpus and is the
core structural idea repeatedly relied on.

\paragraph{Lesson 5: Non-distributivity is not a bug.}
A persistent concern about the nullification operator is its
non-distributivity: $\N(\Delta, x + y) \neq \N(\Delta, x) + \N(\Delta, y)$
in general.
The formalization clarifies that this is a \emph{feature}: it captures
precisely the threshold-crossing behaviour responsible for the spectral gap.
Conservation ($\N + R = \mathrm{id}$) holds exactly throughout; what fails
is linearity, and that failure is the source of the gap.

%% ─────────────────────────────────────────────────────────────────────────────
\section*{Reproducibility}
\label{sec:reproducibility}

The full development compiles under the following pinned environment.

\begin{center}
\renewcommand{\arraystretch}{1.25}
\begin{tabular}{ll}
  \toprule
  \textbf{Component}   & \textbf{Version / Identifier} \\
  \midrule
  Lean                 & \texttt{leanprover/lean4:v4.24.0} \\
  Mathlib commit       & \texttt{f897ebcf72cd16f89ab4577d0c826cd14afaafc7} \\
  Aristotle ATP        & \url{https://aristotle.harmonic.fun} \\
  GitHub repository    & \url{https://github.com/Joeyxyxyz/messina-nullification-framework} \\
  Zenodo Vol.~1 (Core) & \href{https://doi.org/10.5281/zenodo.18902403}{10.5281/zenodo.18902403} \\
  Zenodo Vol.~2 (RH)   & \href{https://doi.org/10.5281/zenodo.18902325}{10.5281/zenodo.18902325} \\
  \bottomrule
\end{tabular}
\end{center}

To build from scratch:
\begin{lstlisting}[language=Lean4]
-- In the repository root:
-- lake update
-- lake build
-- All 82 files should compile with zero sorry diagnostics.
\end{lstlisting}

The repository includes a \texttt{lakefile.lean} pinning the Mathlib commit above.
All proofs relevant to the main theorem (Sections~\ref{sec:proved}--\ref{sec:obstruction})
are zero-\lean{sorry} and zero-\lean{admit}.
The \lean{HPContext} identification hypotheses are the only named assumptions not
discharged by \Mathlib{}; they are explicitly marked in the source.

\subsection*{Independent kernel verification}
\label{sec:axle}

The load-bearing algebraic content of every paper section has been
independently kernel-verified by AXLE (\emph{Axiom Lean Engine},
\url{https://axle.axiommath.ai}, lean-4.28.0 environment).
For each paper section, a self-contained pair of Lean files
(\texttt{*\_STATEMENT.lean}, \texttt{*\_CONTENT.lean}) was submitted to
AXLE's \texttt{verify\_proof} endpoint.
AXLE confirmed each pair: every theorem stated with \lean{:= by sorry}
in the formal-statement side is implemented with a non-\lean{sorry}
proof on the content side, and the implementation type-checks against
the statement.

\begin{center}
\renewcommand{\arraystretch}{1.3}
\footnotesize
\begin{tabular}{lll}
  \toprule
  \textbf{Section} & \textbf{Theorem(s)} & \textbf{AXLE request ID} \\
  \midrule
  \S\ref{sec:operator} & \lean{master\_enforcement} & \texttt{e0755420-c907-4578-8f2f-ab59d002ed52} \\
  \S\ref{sec:proved}   & \lean{sturm\_liouville\_simple\_eigenvalues} (architecture) & \texttt{da53f7a2-78f9-446d-9dda-a130631e6994} \\
  \S\ref{sec:leveling} & \lean{leveling\_theorem}, \lean{leveling\_implies\_re\_half} & \texttt{376a6860-4bb2-4d3d-a461-cce659917ee6} \\
  \S\ref{sec:minimal}  & \lean{resolution\_from\_three\_axioms} (and corollaries) & \texttt{5ac337c1-6c9d-4390-bcaa-42c626f75582} \\
  \bottomrule
\end{tabular}
\end{center}

An earlier audit batch (March 5, 2026, AXLE launch day) verified six
related theorems from the operator-algebra and shifted-spectrum layers
of the corpus:
\texttt{22abc227}, \texttt{572b70d1}, \texttt{c133ead0},
\texttt{3de41f1f}, \texttt{f854b8e9}, \texttt{01ae8bf9}.

The four \texttt{*\_STATEMENT.lean}\,/\,\texttt{*\_CONTENT.lean} pairs
are included in the Zenodo deposit (folder \texttt{axle\_pairs/}) so
reviewers may re-run the verifications independently.
The pairs are designed for AXLE's single-call interface and do not
require a local Lean toolchain; pasting each pair into the AXLE web
console reproduces the verification.

\begin{remark}
The AXLE pairs verify the algebraic logic of each section in
isolation: they expose the heavyweight analytic ingredients (e.g.\
$\mathrm{order}_\rho(\zeta) \geq 1$, the $L^2$ Wronskian decay) as
named hypotheses rather than re-running their full Mathlib-backed
proofs.
This is the appropriate granularity for an independent algebraic
audit.
The full proofs of the named ingredients live in the corpus
(\lean{MinimalAxiomReduction.lean}, \lean{SturmLiouvilleSimplicity.lean})
and are confirmed by the full \texttt{lake build} against the pinned
Mathlib commit.
\end{remark}

%% ─────────────────────────────────────────────────────────────────────────────
\section{Conclusion and Future Directions}
\label{sec:conclusion}

We have presented two Lean~4/Mathlib contributions.
First, a standalone Sturm--Liouville formalization proving eigenvalue simplicity
for confining potentials, independent of any number-theoretic application.
Second, a modular zeta-zero assumption interface expressing RH and Simple Zeros
as conditional consequences of a minimal, explicitly named hypothesis pair,
with all remaining classical ingredients supplied by Mathlib.
The leveling trick provides a novel constructive mechanism showing that a
spectral gap forces non-trivial zeros onto the critical line.
The single remaining open identification step --- connecting the operator's
derived gap to the actual zeta-zero spectrum --- is isolated as an explicit
named context field \lean{hp\_correspondence}, enabling future Mathlib
developments to discharge it.

\paragraph{Future directions.}
\begin{enumerate}
  \item \textbf{Closing SpectralGap.}
        The Hilbert--P\'olya program --- constructing a self-adjoint operator
        whose eigenvalues are the non-trivial zeros --- would make SpectralGap
        a theorem.
        The Sturm--Liouville formalization in this framework (Outputs 59, 62, 64--65)
        is a step toward this goal.

  \item \textbf{Physical applications.}
        The Sturm--Liouville infrastructure developed here is intended as a
        formal bridge toward the Yang--Mills mass gap problem.
        A separate paper formalizing the mass gap approach is in preparation.

  \item \textbf{Mathlib integration.}
        Key infrastructure theorems (the leveling trick, spectral gap property,
        trinary classification) are candidates for Mathlib contribution,
        making them available to the broader formalized mathematics community.
        The goal of this framework is to make the remaining analytic assumptions
        explicit and modular, enabling future formal developments to replace them
        with proved theorems as \Mathlib{}'s analytic infrastructure grows.
        Each named context field in \lean{HPContext} represents a precise target
        for future Mathlib development.
\end{enumerate}

%% ─────────────────────────────────────────────────────────────────────────────
\section*{Artifact Availability}

The complete Lean~4 source code and documentation are publicly available.
An experimental companion artifact (IBM quantum circuits, replication guide)
is included in the Zenodo deposit but is not required for any formal result.

\begin{itemize}
  \item \textbf{Zenodo Vol.~2 (RH):}
        \url{https://doi.org/10.5281/zenodo.18902325}
  \item \textbf{Zenodo Vol.~1 (Core):}
        \url{https://doi.org/10.5281/zenodo.18902403}
  \item \textbf{GitHub:}
        \url{https://github.com/Joeyxyxyz/messina-nullification-framework}
\end{itemize}

The artifact compiles under \texttt{leanprover/lean4:v4.24.0}
with Mathlib commit \texttt{f897ebcf72cd16f89ab4577d0c826cd14afaafc7}.

%% ─────────────────────────────────────────────────────────────────────────────
\section*{Acknowledgements}

The author thanks the Lean community and Mathlib contributors for the
zeta function infrastructure that made this formalization possible,
and the IBM Quantum team for hardware access.
Special thanks to the \Aristotle{} team at Harmonic for automated proof
assistance.

%% ─────────────────────────────────────────────────────────────────────────────
\appendix

\section{Motivational Comparisons and External Interpretations}
\label{app:motivation}

This appendix collects motivational material that informs the design of
$\N(\Delta, x)$ but is not part of the formal contribution of this paper.
None of the statements below are required by any theorem in
Sections~\ref{sec:operator}--\ref{sec:obstruction}; they are included
solely for readers who want broader context.

\subsection{Routing vs.\ Forgetting: $\N(\Delta,x)$ and Robinson's $\mathrm{st}()$}
\label{app:routing}

A natural point of comparison for the nullification operator is
Robinson's \emph{standard part function} $\mathrm{st}(x)$ from
non-standard analysis (NSA)~\cite{robinson1966}.
For a finite hyperreal $x = r + \varepsilon$ where $r$ is real and
$\varepsilon$ is infinitesimal, $\mathrm{st}(r + \varepsilon) = r$.
The infinitesimal part is discarded; $\mathrm{st}()$ is a ring
homomorphism and a quotient map, identifying everything that differs by
an infinitesimal.
We sketch here how $\N(\Delta, x)$ differs from $\mathrm{st}()$
algebraically.

\paragraph{The foundational distinction.}
$\mathrm{st}()$ is a \emph{forgetting} operation: sub-threshold
(infinitesimal) content is discarded entirely.
$\N(\Delta, x)$ is a \emph{routing} operation: sub-threshold content is
sent to a separate channel via the companion reservoir operator
$R(\Delta, x) = x - \N(\Delta, x)$, where it remains accessible and
recoverable via $R$.
The conservation law $\N(\Delta, x) + R(\Delta, x) = x$ holds for all
$x$ and all $\Delta > 0$.
There is no analog of this identity in NSA: $\mathrm{st}()$ has no
conservation partner of this form.

\paragraph{Trinary classification.}
$\N(\Delta, x)$ generates a three-class partition of $\R$:
\begin{itemize}[noitemsep]
  \item $[0]$ --- true zero ($x = 0$);
  \item $[0.0e]$ --- nullified/reservoir ($0 < |x| < \Delta$);
  \item $[1]$ --- observable ($|x| \geq \Delta$).
\end{itemize}
The middle class is the structural innovation.
It is absent from frameworks built around $\mathrm{st}()$, which
collapses all infinitesimals to zero, and from standard binary threshold
models, which have no mechanism to retain sub-threshold magnitude.

\paragraph{Non-distributivity.}
$\mathrm{st}()$ is distributive ($\mathrm{st}(x+y) = \mathrm{st}(x) +
\mathrm{st}(y)$).
$\N(\Delta, x)$ is not: in general
$\N(\Delta, x + y) \neq \N(\Delta, x) + \N(\Delta, y)$.
Two sub-threshold components can sum to a super-threshold result.
This is not a defect; it is precisely the algebraic content of
threshold-crossing.

\paragraph{Comparison summary.}
\begin{center}
\renewcommand{\arraystretch}{1.3}
\begin{tabular}{lll}
  \toprule
  \textbf{Property} & \textbf{Robinson $\mathrm{st}()$} & \textbf{$\N(\Delta, x)$} \\
  \midrule
  Type              & Ring homomorphism (quotient)   & Filtration operator (routing) \\
  Sub-threshold     & Collapsed to zero       & Routed to reservoir $[0.0e]$ \\
  Information       & Discarded               & Magnitude and sign retained \\
  Distributive?     & Yes                     & No (structural, not incidental) \\
  Zero structure    & One class: $\{0\}$      & Three classes: $[0]$, $[0.0e]$, $[1]$ \\
  Signed extension  & Not defined             & $[{+}0.0e]$, $[{-}0.0e]$ distinct \\
  \bottomrule
\end{tabular}
\end{center}

For the Lean~4 proofs of non-distributivity, conservation, and the
signed classification system, see
\lean{Batch3\_Distributivity\_aristotle.lean},
\lean{batch11\_nsa\_equivalence\_aristotle.lean}, and
\lean{output71\_signed\_resolution.lean} in the corpus.

\subsection{Concrete Numerical Example of Non-Distributivity}
\label{app:xi}

Non-distributivity has a concrete numerical instantiation that can be
machine-checked.
With threshold $\Lambda = 246$ and values $m_s = 96$, $m_\Xi = 1315$
(formalized in \lean{Batch4.lean}):
\[
  \N(\Lambda,\, m_s) + \N(\Lambda,\, m_s) = 0 + 0 = 0
  \;\neq\; 1315 = \N(\Lambda,\, m_\Xi).
\]
The threshold-crossing behavior --- two sub-threshold components summing to
a super-threshold value --- is precisely the algebraic structure that makes
$\N$ a filtration operator rather than a ring homomorphism.
This is formalized in \lean{Batch4.lean}:

\begin{lstlisting}[language=Lean4]
def threshold : ℕ := 246
def m_component : ℕ := 96    -- below threshold: N(246, 96) = 0
def m_composite : ℕ := 1315  -- above threshold: N(246, 1315) = 1315

theorem non_distributivity_concrete :
    N threshold m_component + N threshold m_component = 0 ∧
    N threshold m_composite = m_composite := by
  simp [N, threshold, m_component, m_composite]
\end{lstlisting}

\begin{remark}
The numerical values above correspond to a known particle physics threshold
(the QCD confinement scale $\Lambda_\text{QCD} \approx 246$\,MeV and
the constituent strange-quark / Xi-baryon masses), independently confirmed
by PDG data~\cite{pdg2024} to sub-MeV precision.
This correspondence is noted as external motivation only; the Lean
theorems above hold for these numerical values regardless of any physical
interpretation.
\end{remark}

\section{Experimental Companion Artifact (non-deductive motivation)}
\label{app:experimental}

This appendix describes a physical experiment that motivates the
\lean{MultipleZeroKillsGap} hypothesis used in Path~C of
Section~\ref{sec:minimal}.
It does not constitute a formal proof and is not required for any theorem
in Sections~\ref{sec:proved}--\ref{sec:obstruction}.
It is included as a reproducible companion artifact for researchers
interested in the empirical instantiation of the framework's spectral
structure.

The \lean{MultipleZeroKillsGap} hypothesis --- that double-zero structure
measurably destroys the spectral gap signature --- was tested on IBM quantum
hardware via 64-qubit lattice gauge theory circuits.

\paragraph{Experiment design.}
Two circuits were run in matched sessions:
\begin{itemize}[noitemsep]
  \item \textbf{C0b} (depth-matched): identity CNOT pairs, same circuit
        depth, zero degeneracy. Noise control.
  \item \textbf{C3} (degeneracy): state-copying CNOTs forcing double-zero
        structure.
\end{itemize}
Both circuits use 64 qubits encoding the lattice Hamiltonian.
Sessions were matched on backend and time to control for drift.

\paragraph{Results.}
\begin{center}
\begin{tabular}{ll}
  \toprule
  Metric & Value \\
  \midrule
  $P(\text{C0b free} > \text{C3 free})$ & $1.0000$ \\
  Bootstrap resamples (5{,}000) & 5{,}000 / 5{,}000 \\
  95\% CI of difference & $[0.011, 0.025]$ \\
  Total measurements & 161{,}220+ \\
  Backends & ibm\_brisbane, ibm\_torino, ibm\_marrakesh \\
  \bottomrule
\end{tabular}
\end{center}

Degeneracy (forced double-zero structure) measurably and reproducibly
suppresses the confinement gap signature relative to matched controls.
These results provide empirical motivation for \lean{MultipleZeroKillsGap}
as a plausible premise for Path~C; they do not constitute a formal proof
of that hypothesis or its analytic consequences.
The full replication guide and circuit files are included in the Zenodo
artifact.

%% ─────────────────────────────────────────────────────────────────────────────
\bibliographystyle{plain}
\begin{thebibliography}{99}

\bibitem{loeffler2025}
D.~Loeffler and M.~Stoll.
\newblock Formalizing zeta and $L$-functions in {Lean}.
\newblock \emph{Annals of Formalized Mathematics}, 2025.
\newblock \href{https://doi.org/10.46298/afm.15328}{DOI:10.46298/afm.15328}.

\bibitem{mathlib2020}
The~Mathlib Community.
\newblock The {Lean} mathematical library.
\newblock In \emph{Proceedings of the 9th ACM SIGPLAN International
  Conference on Certified Programs and Proofs (CPP 2020)}, pages 367--381, 2020.
\newblock \href{https://doi.org/10.1145/3372885.3373824}{DOI:10.1145/3372885.3373824}.

\bibitem{aristotle2025}
Harmonic.
\newblock Aristotle: {AI}-assisted theorem proving in {Lean~4}.
\newblock \url{https://aristotle.harmonic.fun}, 2025.

\bibitem{montgomery1973}
H.~L. Montgomery.
\newblock The pair correlation of zeros of the zeta function.
\newblock In \emph{Analytic Number Theory (Proc. Sympos. Pure Math., Vol. XXIV)},
  pages 181--193. Amer. Math. Soc., 1973.

\bibitem{speiser1935}
A.~Speiser.
\newblock Geometrisches zur {Riemannschen} {Zetafunktion}.
\newblock \emph{Math. Ann.}, 110(1):514--521, 1935.

\bibitem{goldston2001}
D.~A. Goldston, S.~M. Gonek, and H.~L. Montgomery.
\newblock Mean values of the logarithmic derivative of the {Riemann} zeta
  function with applications to primes in short intervals.
\newblock \emph{J. Reine Angew. Math.}, 537:105--126, 2001.

\bibitem{trudgian2014}
T.~S. Trudgian and D.~J. Platt.
\newblock Improved explicit bounds for some functions of prime numbers.
\newblock \emph{J. Number Theory}, 2014.

\bibitem{debruijn1950}
N.~G. de~Bruijn.
\newblock The roots of trigonometric integrals.
\newblock \emph{Duke Math. J.}, 17:197--226, 1950.

\bibitem{moura2021}
L.~de~Moura and S.~Ullrich.
\newblock The {Lean~4} theorem prover and programming language.
\newblock In \emph{Automated Deduction -- CADE 28}, Lecture Notes in
  Computer Science, pages 625--635. Springer, 2021.

\bibitem{robinson1966}
A.~Robinson.
\newblock \emph{Non-Standard Analysis}.
\newblock North-Holland, Amsterdam, 1966.

\bibitem{pdg2024}
Particle Data Group; S.~Navas et al.
\newblock Review of particle physics.
\newblock \emph{Phys.\ Rev.\ D}, 110:030001, 2024.
\newblock \href{https://doi.org/10.1103/PhysRevD.110.030001}{DOI:10.1103/PhysRevD.110.030001}.

\end{thebibliography}

\end{document}
